\documentclass[12pt,oneside]{book}

% ============================================================
% METADATOS
% ============================================================
\newcommand{\universidad}{Universidad de Buenos Aires (UBA)}
\newcommand{\materia}{Lectocomprensión}
\newcommand{\titulo}{Inglés jurídico: el proceso penal en Estados Unidos}
\newcommand{\autor}{Isaac Edgar Camacho Ocampo}
\newcommand{\anio}{Octubre de 2026}

% ============================================================
% PAQUETES
% ============================================================
\usepackage[utf8]{inputenc}
\usepackage[T1]{fontenc}
% Babel no se carga aquí para evitar errores de idioma en instalaciones mínimas de TeX.

\usepackage[
    a4paper,
    top=2.5cm,
    bottom=2.5cm,
    left=3cm,
    right=2.5cm,
    headheight=15pt
]{geometry}

\usepackage{amsmath}
\usepackage{amssymb}
\usepackage{mathtools}

\usepackage{graphicx}
\usepackage{float}
\usepackage{caption}
\usepackage{subcaption}

\usepackage{booktabs}
\usepackage{array}
\usepackage{longtable}
\usepackage{tabularx}

\usepackage{enumitem}

\usepackage{xcolor}
\usepackage[most]{tcolorbox}

\usepackage{fancyhdr}
\usepackage{ragged2e}

\usepackage[
    colorlinks=true,
    linkcolor=blue!50!black,
    citecolor=green!40!black,
    urlcolor=blue!60!black,
    pdftitle={\titulo},
    pdfauthor={\autor},
    pdfsubject={Apuntes universitarios},
    bookmarks=true
]{hyperref}

% ============================================================
% CONFIGURACIÓN
% ============================================================
\graphicspath{
    {./img/}
    {./graficos/}
    {./figuras/}
}

\setcounter{tocdepth}{2}

\setlength{\parindent}{0pt}
\setlength{\parskip}{0.5em}

\setlist[itemize]{
    topsep=0.3em,
    itemsep=0.2em
}

\setlist[enumerate]{
    topsep=0.3em,
    itemsep=0.2em
}

\clubpenalty=10000
\widowpenalty=10000

\captionsetup{font=small,labelfont=bf}

\newcolumntype{L}{>{\raggedright\arraybackslash}X}

% ============================================================
% COMANDOS
% ============================================================
\newcommand{\abs}[1]{\left|#1\right|}
\newcommand{\norm}[1]{\left\lVert#1\right\rVert}
\newcommand{\vect}[1]{\mathbf{#1}}
\newcommand{\deriv}[2]{\frac{d#1}{d#2}}
\newcommand{\pderiv}[2]{\frac{\partial#1}{\partial#2}}

% ============================================================
% ENTORNOS / CAJAS
% ============================================================
\newtcolorbox{definition}[1]{
    enhanced,
    breakable,
    title={Definición: #1},
    fonttitle=\bfseries,
    colback=blue!3,
    colframe=blue!50!black,
    boxrule=0.8pt,
    arc=2mm
}

\newtcolorbox{important}{
    enhanced,
    breakable,
    title={Importante},
    fonttitle=\bfseries,
    colback=yellow!8,
    colframe=orange!70!black,
    boxrule=0.8pt,
    arc=2mm
}

\newtcolorbox{example}[1]{
    enhanced,
    breakable,
    title={Ejemplo: #1},
    fonttitle=\bfseries,
    colback=green!4,
    colframe=green!40!black,
    boxrule=0.8pt,
    arc=2mm
}

\newtcolorbox{formula}[1]{
    enhanced,
    breakable,
    title={Fórmula / Regla: #1},
    fonttitle=\bfseries,
    colback=purple!4,
    colframe=purple!50!black,
    boxrule=0.8pt,
    arc=2mm
}

\newtcolorbox{exercise}[1]{
    enhanced,
    breakable,
    title={Ejercicio: #1},
    fonttitle=\bfseries,
    colback=gray!5,
    colframe=gray!60!black,
    boxrule=0.8pt,
    arc=2mm
}

\newtcolorbox{note}{
    enhanced,
    breakable,
    title={Nota},
    fonttitle=\bfseries,
    colback=white,
    colframe=gray!50,
    boxrule=0.6pt,
    arc=2mm
}

% ============================================================
% CONFIGURACIÓN FINAL
% ============================================================
\pagestyle{fancy}
\fancyhf{}

\fancyhead[L]{\nouppercase{\leftmark}}
\fancyhead[R]{\materia}

\fancyfoot[C]{\thepage}

\renewcommand{\headrulewidth}{0.4pt}
\renewcommand{\footrulewidth}{0pt}

\fancypagestyle{plain}{
    \fancyhf{}
    \fancyfoot[C]{\thepage}
    \renewcommand{\headrulewidth}{0pt}
}

% ============================================================
% DOCUMENTO
% ============================================================
\begin{document}

% ------------------------------------------------------------
% PORTADA
% ------------------------------------------------------------
\begin{titlepage}

\thispagestyle{empty}

\begin{center}

\vspace*{1cm}

{\large \textsc{\universidad}}

\vspace{3cm}

{\Huge \textbf{\materia}}

\vspace{1cm}

{\LARGE \titulo}

\vspace{0.8cm}

\small\textit{Leer con atención es el primer paso para convertir el estudio en conocimiento.}
\normalsize

\vspace{1.2cm}

\rule{0.70\textwidth}{0.5pt}

\vspace{0.5cm}

{\large Apuntes de estudio}

\vspace{0.5cm}

\rule{0.70\textwidth}{0.5pt}

\vfill

{\large \autor}

\vspace{0.3cm}

{\large \anio}

\end{center}

\vfill

\begin{flushleft}
\textit{Este es un modesto aporte a los estudiantes de ``Lectocomprensión'' no pretende sustituir el material oficial de la materia.}
\end{flushleft}

\end{titlepage}

% ------------------------------------------------------------
% ÍNDICE
% ------------------------------------------------------------
\frontmatter
\tableofcontents
\mainmatter

% ------------------------------------------------------------
% CONTENIDO
% ------------------------------------------------------------

% ============================================================
\chapter[Presentación de la clase]{Presentación de la clase}
% ============================================================

\section{Contenido general}

La clase recorre, a partir de un texto en inglés, las etapas del proceso penal de Estados Unidos: gran jurado, audiencias previas, mociones, juicio ante jurado, veredicto, \textit{hung jury}, \textit{mistrial}, determinación de la pena y \textit{allocution}. Luego analiza el vocabulario y las garantías del debido proceso en las Enmiendas Quinta y Sexta, la relación entre el \textit{U.S. Code} y los códigos estatales, y las reglas de procedimiento policial (Miranda).

\section{Utilidad profesional de la materia}

La clase destaca cinco aportes de la lectocomprensión en inglés jurídico a la vida profesional:

\begin{itemize}
    \item \textbf{Leer fuentes reales sin depender de traducciones.} Fallos, contratos y normas internacionales están en inglés. Entender términos como \textit{indictment}, \textit{motion} o \textit{due process} permite trabajar con la fuente original.
    \item \textbf{Comparar sistemas jurídicos.} Advertir que la \textit{probation} estadounidense es una pena y que la figura argentina es una suspensión del juicio a prueba evita errores graves de traducción y de concepto.
    \item \textbf{Pensar el propio derecho con mirada crítica.} Contrastar estándar probatorio, jurado, unanimidad o plazos ayuda a argumentar mejor sobre garantías y reformas.
    \item \textbf{Evitar los ``falsos amigos''.} Palabras como \textit{sentence} o \textit{property} parecen conocidas y significan otra cosa; en un escrito o una negociación, el error es profesional.
    \item \textbf{Entrenar una habilidad transferible.} Identificar términos resaltados, anticipar el tema y leer sin frenarse en cada palabra sirve para cualquier texto técnico, en cualquier idioma.
\end{itemize}

\section{Organización de los contenidos}

Los contenidos se agrupan en cinco bloques que guardan relación entre sí:

\begin{enumerate}
    \item \textbf{Etapas previas al juicio:} gran jurado y audiencia preliminar, imputación y procesamiento, y mociones previas al juicio.
    \item \textbf{El juicio y el jurado:} etapa de debate, los tres jurados, roles del juez y del jurado, instrucciones, deliberación, veredicto, \textit{hung jury} y \textit{mistrial}.
    \item \textbf{Sentencia y pena:} penas posibles, \textit{sentencing jury} y \textit{allocution}.
    \item \textbf{Garantías constitucionales:} debido proceso, Quinta Enmienda y Sexta Enmienda.
    \item \textbf{Normas y procedimiento policial:} \textit{U.S. Code} y códigos estatales, \textit{Federal Rules} como piso mínimo, detención policial y Miranda.
\end{enumerate}

\begin{note}
La clase propuso una analogía con la construcción de una casa: primero se revisan permisos y materiales (si un material no cumple las normas, se descarta y la obra puede caerse; es el bloque 1); luego se levanta la obra a la vista de todos, con inspectores que deben coincidir en que está bien hecha (bloque 2: el juicio y el jurado); después llega la habilitación final, que fija qué se aprueba o se corrige (bloque 3: la pena). Todo el tiempo rige un reglamento de construcción que nadie puede violar (bloque 4: las garantías constitucionales) y existen códigos locales que varían de una ciudad a otra, además de pautas para quienes controlan la obra (bloque 5: códigos federales y estatales, y procedimiento policial).
\end{note}

% ============================================================
\chapter[Etapas previas al juicio]{Etapas previas al juicio}
% ============================================================

\section{Investigación y gran jurado}

Lo estudiado hasta este punto (\textit{grand jury} y \textit{preliminary hearing}) corresponde a lo que en el sistema argentino se denomina etapa de instrucción o investigación, que en nuestro sistema puede llevar adelante el fiscal a nivel judicial. En el sistema estadounidense interviene además el gran jurado cuando se trata de delitos aberrantes o muy severos, o de delitos a los que puede corresponder la pena de muerte: delitos graves y a nivel federal.

Tras la audiencia preliminar comienza una etapa vinculada ya con lo probatorio. Se trata de dos audiencias distintas. En este punto está, en términos del sistema argentino, trabada la litis: la persona conoce completamente de qué se la acusa y se ha decidido continuar con el proceso.

\section{Imputación y procesamiento}

El docente señaló que suele confundir ambas nociones y recurrió a los estudiantes que ya cursaron Derecho Penal para precisar el orden. Lo establecido en clase es lo siguiente: primero se imputa y, cuando sucede la acusación formal, la persona queda procesada; para ello tiene que haber una cierta certeza y más pruebas.

Con esa base, y aunque las figuras no sean idénticas, puede forzarse una asimilación: después del \textit{indictment}, o de la etapa del gran jurado, hay un procesamiento. Como lo entendemos nosotros, el procesamiento supone que ya se evaluó que existen indicios probatorios suficientes y que se decide continuar con el proceso.

\begin{table}[H]
\centering
\small
\caption{Equivalencia aproximada de etapas (según lo expuesto en clase)}
\begin{tabularx}{\textwidth}{L L L}
\toprule
\textbf{Sistema argentino} & \textbf{Sistema de Estados Unidos} & \textbf{Qué se evalúa} \\
\midrule
\textbf{Imputación} & \textit{Arrest} / \textit{Preliminary hearing} & Existencia de elementos para sostener una acusación inicial. \\
\addlinespace
\textbf{Procesamiento} & \textit{Indictment} (acusación formal, tras el gran jurado) o \textit{information} & Indicios probatorios suficientes para continuar con el proceso. \\
\addlinespace
\textbf{Etapa probatoria / debate} & \textit{Trial} (precedido de \textit{pretrial motions}) & Producción de prueba ante el jurado y veredicto. \\
\bottomrule
\end{tabularx}
\end{table}

\section{Mociones previas al juicio (\textit{pretrial motions})}

\begin{definition}{Motion}
Petición, pedido o solicitud de una de las partes al juez para que resuelva o haga algo. Es una categoría amplia, que también se vio en el proceso civil.
\end{definition}

En esta instancia las \textit{motions} se vinculan con la prueba. El término \textit{pretrial} no alude a algo perjudicial: significa que son previas a la etapa probatoria, cuando ya se procesó a la persona. El texto de la clase menciona, por ejemplo, \textit{to suppress evidence}: suprimir o no admitir ciertos medios de prueba, entre ellos la prueba testimonial. También pueden solicitarse desestimaciones de la acusación por insuficiencia de pruebas o por falta de sustento de las imputaciones. En definitiva, son distintos pedidos al juez para que resuelva cuestiones procesales, y el juez puede hacer lugar o no.

\subsection{Por qué se revisa nuevamente la prueba}

Surgió la objeción de que resulta extraño volver a hablar de prueba cuando ya se la había evaluado para la acusación formal. La explicación dada es que se trata de una segunda revisión: antes del procesamiento se determinó que había indicios probatorios suficientes para avanzar, pero al llegar a la audiencia previa al juicio la contraparte puede alegar, por ejemplo, que determinada prueba fue obtenida ilícitamente o que un testigo es amigo, primo o hermano de alguien. Al suprimirse evidencia puede suceder que la prueba restante sea insuficiente; entonces corresponde pedir la desestimación de la acusación, aunque ya hubiera sido evaluada antes.

\begin{example}{Cadena de custodia}
Un análisis de alcoholemia, un análisis psicológico o el estudio de un cuchillo cuya cadena de custodia no fue bien tramitada. La defensa objeta cómo se trató esa prueba y, si la prueba se cae, y era lo más sólido que tenía la parte acusadora, la acusación se debilita. La prueba queda viciada.
\end{example}

Por eso es tan importante la etapa del arresto y el cumplimiento de las garantías: puede ser claro lo que sucedió y que la persona sea efectivamente culpable, pero si procesalmente se impugnan las pruebas y se caen, los hechos son los hechos, pero no hay prueba.

Sobre la frecuencia de esta situación, el docente manifestó no creer que suceda con mucha frecuencia y reconoció que lo desconoce; indicó que puede pasar que las pruebas se vayan declarando inadmisibles.

\begin{important}
Las \textit{pretrial motions} son pedidos al juez, previos a la etapa probatoria, para suprimir pruebas (por ejemplo, las obtenidas ilícitamente o con cadena de custodia viciada) o incluso para desestimar la acusación. Por eso es tan importante que el arresto y la recolección de prueba respeten las garantías: una prueba contundente puede caerse por defectos de procedimiento.
\end{important}

\section{Cuadro Cornell: etapas previas al juicio}

\begin{longtable}{@{}>{\raggedright\arraybackslash}p{0.28\textwidth}>{\raggedright\arraybackslash}p{0.66\textwidth}@{}}
\toprule
\textbf{Preguntas / palabras clave} & \textbf{Notas / respuestas} \\
\midrule
\endfirsthead
\toprule
\textbf{Preguntas / palabras clave} & \textbf{Notas / respuestas} \\
\midrule
\endhead
\bottomrule
\endlastfoot

\textbf{¿Cuándo interviene el gran jurado y qué decide?} &
Interviene solo en delitos federales muy graves o con posible pena de muerte. Decide si hay indicios suficientes para avanzar con la acusación (\textit{indictment}); cumplida su función, se lo despide. \\
\addlinespace

\textbf{¿Qué va primero: imputación o procesamiento?} &
Primero la imputación; luego, con la acusación formal y tras evaluar indicios probatorios suficientes, el procesamiento. El \textit{indictment} se asimila aproximadamente al procesamiento. \\
\addlinespace

\textbf{¿Qué es una \textit{pretrial motion}?} &
Petición al juez, previa a la etapa probatoria, para suprimir pruebas (ilícitas, con cadena de custodia viciada, testigos cuestionados) o desestimar la acusación. Funciona como una segunda revisión de la prueba. \\
\addlinespace

\textbf{¿Por qué es tan importante la etapa del arresto?} &
Porque si el procedimiento no respeta las garantías, la prueba puede impugnarse y caerse, aun cuando los hechos sean claros. \\

\midrule
\multicolumn{2}{@{}p{\dimexpr0.94\textwidth+2\tabcolsep\relax}@{}}{%
\textbf{Resumen:} Tras el gran jurado y el \textit{indictment} (asimilable al procesamiento), las \textit{pretrial motions} permiten depurar la prueba antes del juicio y hasta lograr la desestimación de la acusación. Por eso el respeto de las garantías desde el arresto es decisivo.} \\
\end{longtable}

% ============================================================
\chapter[El juicio y el jurado]{El juicio (\textit{trial}) y el jurado}
% ============================================================

\section{La etapa de debate: apertura, prueba y cierre}

Una vez resueltos los pedidos previos (el juez puede admitirlos o no), comienza la etapa que, al igual que en el proceso civil, se denomina \textit{trial}: la etapa probatoria, que en nuestro sistema corresponde a la etapa del debate.

Allí sucede lo mismo que en el proceso civil: los \textit{opening statements}, toda la producción de prueba (medios de prueba y declaraciones) y luego el \textit{closing statement} o \textit{closing argument}. Son los alegatos de apertura y de cierre, con 15 minutos para cada parte; comienza la fiscalía.

\subsection{El jurado en el juicio}

En el proceso penal el jurado se compone de 12 personas. En general se convoca a más de 12 (el docente mencionó alrededor de 20) porque debe realizarse el proceso de selección de los 12 jurados, y se convocan algunos suplentes por si ocurre algún imprevisto.

Los alegatos de apertura se hacen ya ante el jurado. Cada parte anticipa su posición: la acusación sostiene que probará que la persona cometió el delito y cómo ocurrió; la defensa afirma lo contrario y, si hay una coartada, la plantea, aunque sea muy brevemente. Luego se produce toda la prueba y el cierre sintetiza lo producido.

\subsection{Orden de intervención de las partes}

Según el texto, la defensa puede hacer su alegato de apertura hasta que la fiscalía termine su presentación (\textit{until the prosecution rests}). Esto permite a la defensa escuchar primero lo que se dice en su contra y reservarse la última palabra, de modo que lo último que escuche el jurado sea la defensa. Primero presenta sus testigos y pruebas la fiscalía; después lo hace la defensa, y cuando esta termina su presentación se dice que la defensa ``descansa'' (\textit{defense rests}). En los alegatos de cierre el orden es inverso: primero la defensa y después la fiscalía.

\section{Los tres jurados: \textit{grand jury}, \textit{petit jury} y \textit{sentencing jury}}

El \textit{indictment} no lo dicta el mismo jurado que interviene en el juicio. El \textit{grand jury} es otro jurado y solo interviene, como se explicó al comienzo, en causas muy graves: delitos federales muy graves y delitos graves en los que la pena puede ser la pena de muerte; en esos casos se convoca siempre un gran jurado. En el juicio interviene el jurado ordinario, que evalúa la prueba.

\begin{definition}{Grand jury (gran jurado)}
Jurado que interviene antes del juicio, solo en delitos federales muy graves o con posible pena de muerte. Su función es determinar si hay suficiente para avanzar con la acusación, es decir, si existen indicios suficientes para pensar que la persona es responsable de lo que se le atribuye. Si responde afirmativamente, se avanza y el jurado es despedido; no continúa hasta el final del proceso. Se convoca entonces al jurado del juicio.
\end{definition}

En una misma causa puede llegar a haber tres jurados, como se muestra en el siguiente cuadro.

\begin{table}[H]
\centering
\small
\caption{Los tres jurados del proceso penal estadounidense}
\begin{tabularx}{\textwidth}{L L L}
\toprule
\textbf{Jurado} & \textbf{Cuándo interviene} & \textbf{Qué decide} \\
\midrule
\textbf{\textit{Grand jury} (gran jurado)} & Antes del juicio; solo en delitos federales muy graves o con posible pena de muerte. & Si hay indicios suficientes para avanzar con la acusación (\textit{indictment}). \\
\addlinespace
\textbf{\textit{Petit jury} (jurado ordinario, 12 personas)} & Durante el \textit{trial}. & Evalúa la prueba y emite el veredicto: \textit{guilty} / \textit{not guilty}. \\
\addlinespace
\textbf{\textit{Sentencing jury} (jurado de sentencia)} & Después del veredicto de culpabilidad; solo en delitos gravísimos. & Qué pena corresponde (ya no se discute la culpabilidad). \\
\bottomrule
\end{tabularx}
\end{table}

\section{Roles del juez y del jurado}

El trabajo del juez es de control y de reglas de juego: dicta las reglas, ordena y supervisa, y resuelve las \textit{motions} cuando determina que una prueba es admisible o no. Juez y jurado trabajan en conjunto, pero con roles bien diferenciados: el jurado evalúa lo fáctico y el juez lo jurídico, porque el jurado no debe tener conocimientos de derecho.

\begin{table}[H]
\centering
\small
\caption{División de roles}
\begin{tabularx}{\textwidth}{L L}
\toprule
\textbf{Juez} & \textbf{Jurado} \\
\midrule
Controla las reglas de juego y supervisa el proceso. & Evalúa los hechos (lo fáctico). \\
\addlinespace
Resuelve las \textit{motions}: decide qué prueba es admisible. & Valora la prueba admitida. \\
\addlinespace
Aporta lo jurídico; impone la pena (en general). & No debe tener conocimientos jurídicos; emite el veredicto. \\
\bottomrule
\end{tabularx}
\end{table}

\section{Costo y duración del proceso con jurado}

El docente observó que le faltan datos reales, como estadísticas o el relato de un abogado que ejerza en Estados Unidos, sobre cuán costoso es, en dinero y en tiempo, un proceso en el que puede producirse la anulación del juicio, y con qué frecuencia ocurre. Cada proceso de este tipo supone tiempo, recursos del juzgado, tiempo de los abogados (que se factura) y toda la selección de jurados. A ello se suma que alguien debe instruir y atender a los jurados, sobre todo si se los aísla. Por sentido común, el montaje es muy costoso, ni hablar en tiempo.

Desde el sistema argentino se suele ver el proceso con otros ojos, con la expectativa de que sea más rápido. El docente señaló que aquí también las cosas podrían ser más rápidas y mencionó los plazos procesales: dentro de diez días se presenta un escrito, dentro de otros diez, y en el medio pasaron tres meses y la causa sigue en el mismo lugar.

\section{Instrucciones al jurado}

Antes de que el jurado delibere, el juez le da instrucciones sobre cómo analizar la prueba, pues durante la producción de la prueba pueden haber ocurrido situaciones que influyan en su apreciación.

\begin{example}{Fallo sobre presión al jurado}
Según un fallo de uno de los textos, que terminó en apelación, se trataba de un homicidio. Los familiares de la víctima asistieron a la audiencia pública con carteles, pancartas, pines y fotos de la víctima, lo que resulta impactante para el jurado. No se recordaba con precisión si había existido una \textit{motion} para pedir al juez que evitara esa situación y este la denegó, o si ni siquiera se presentó. % TODO: contenido incompleto o ambiguo en las notas originales
En la apelación se logró cuestionar el veredicto, con el argumento de que el jurado se había visto muy impactado por lo ocurrido en la audiencia. Se consideró que hubo un error procesal por no haber actuado el juez para impedirlo.
\end{example}

En las instrucciones el juez da las pautas sobre cómo apreciar la prueba escuchada. Cuando una de las partes objeta una declaración y el juez hace lugar a la objeción, la porción del testimonio se quita del expediente: queda registrado que no se tiene como válida y se indica al jurado que no la tenga en cuenta. Pero el jurado ya la escuchó, y es muy difícil que logre ``borrarla'': se trata de personas comunes a quienes se pide, por ejemplo, que no consideren las imágenes de la víctima, y ello toca cuestiones humanas, emociones y sentimientos. El docente consideró interesante investigar cómo se maneja este problema, y subrayó que es esencial que las instrucciones sean muy claras sobre qué deben tener en cuenta los jurados y qué no.

\section{Deliberación, estándar probatorio y veredicto}

Después de las instrucciones, los jurados se retiran a deliberar, igual que en el jurado civil. El estándar probatorio en el proceso penal es \textit{beyond a reasonable doubt}: más allá de toda duda razonable. Los 12 integrantes deben determinar simplemente ``sí'' o ``no''; por decir que no se declara ``inocente'' (\textit{not guilty}). Para declarar culpable a una persona es necesario cumplir con ese estándar.

\begin{formula}{Estándar probatorio penal}
\textit{Beyond a reasonable doubt} (más allá de toda duda razonable). Para el veredicto de culpabilidad se exige unanimidad de los 12 jurados.
\end{formula}

Los jurados votan como quieran (el docente indicó que tiene entendido que lo hacen con un papelito) y a puerta cerrada, sin necesidad de justificar su voto. Si uno solo de los 12 no tiene certeza absoluta de la culpabilidad y vota ``no culpable'', hay que acordar la cuestión y volver a pensarla. Si pasa el tiempo sin acuerdo, lo que terminan por declarar es ``no culpable'': no porque estén convencidos de la inocencia, sino porque no están convencidos de la culpabilidad y no alcanzan el estándar. El veredicto debe ser unánime y se anuncia en conjunto: habla uno solo, pero anuncian que todos juntos llegaron a la decisión.

Durante la deliberación los jurados debaten entre sí respecto de toda la prueba presentada: intentan convencerse, se dan opiniones, y en algún momento votan; si no se ponen de acuerdo, vuelven a discutir. No son llevados a otro lugar para volver a decidir después. Si en un proceso largo no recuerdan, por ejemplo, la declaración de un testigo de hace 20 días, pueden pedir aclaraciones al juez que los ayuden a resolver, y vuelven a deliberar.

\subsection{Comparación con el proceso civil}

\begin{table}[H]
\centering
\small
\caption{Comparación entre el proceso penal y el civil en Estados Unidos}
\begin{tabularx}{\textwidth}{>{\raggedright\arraybackslash}p{0.2\textwidth} L L}
\toprule
\textbf{Aspecto} & \textbf{Proceso penal} & \textbf{Proceso civil} \\
\midrule
\textbf{Estándar probatorio} & \textit{Beyond a reasonable doubt} (más allá de toda duda razonable). & Preponderancia de la prueba (probabilidades); estándar más laxo. \\
\addlinespace
\textbf{Integrantes del jurado} & 12 personas. & 6 personas (ejemplo dado en clase). \\
\addlinespace
\textbf{Unanimidad} & Obligatoria: basta un voto discordante para que no haya condena. & No necesariamente: en el ejemplo, bastan 4 de 6. \\
\addlinespace
\textbf{Veredicto} & \textit{Guilty} / \textit{not guilty} (culpable / inocente), anunciado en conjunto. & Responsabilidad; también anunciado en conjunto. \\
\addlinespace
\textbf{Alegatos} & Apertura y cierre; 15 minutos por parte; comienza la fiscalía. & Exactamente igual: \textit{opening} y \textit{closing statements}. \\
\addlinespace
\textbf{Etapa probatoria} & Producción de prueba, testigos y declaraciones. & Exactamente igual. \\
\bottomrule
\end{tabularx}
\end{table}

En el proceso civil el estándar es la preponderancia de la prueba, una cuestión de probabilidades: se trata de decidir hacia qué lado se inclina el jurado, si piensa que la persona es responsable o no. Es un estándar mucho más laxo y no se requiere necesariamente unanimidad: si, de seis integrantes, dos dudan, con cuatro pueden anunciar el veredicto. En ambos casos el veredicto se anuncia en conjunto; lo ocurrido en la sala de deliberación es un asunto del jurado, salvo que se lo convoque por una duda o un problema.

\begin{important}
Salvo la parte de las distintas audiencias previas, el proceso penal estadounidense es muy parecido al civil. Las diferencias principales son el estándar probatorio (más allá de toda duda razonable) y la cantidad de miembros del jurado (con exigencia de unanimidad).
\end{important}

\subsection{Aislamiento del jurado}

El aislamiento del jurado se da, en general, durante la etapa probatoria, para evitar que lo influyan la opinión mediática, los medios y la opinión pública. Durante la deliberación hay una sala de deliberación contigua a la sala de audiencias; si la deliberación se prolonga, el docente supuso que se los lleva a un hotel para descansar, y señaló que pueden ser retenidos si deliberan durante mucho tiempo.

\section{\textit{Hung jury} y \textit{mistrial}}

Puede ocurrir que el jurado vote una vez, dos veces o más, y no se ponga de acuerdo. Puede acordar un veredicto de \textit{not guilty} aun cuando haya uno solo con dudas. Pero si no logra consensuar qué veredicto emitir, informa al juez que no pudo llegar a un veredicto; puede estar tan trabado que nadie quiera ceder. La división puede darse en muchas combinaciones (por ejemplo, cinco y siete, y no solo un jurado con dudas). El jurado no puede forzar un veredicto: debe decir que no se ponen de acuerdo.

\begin{definition}{\textit{Hung jury}}
Jurado bloqueado o estancado: no logra unanimidad ni para condenar ni para absolver. No importa cuántos discrepen; no es necesariamente seis contra seis (puede ser, por ejemplo, once contra uno).
\end{definition}

\begin{definition}{\textit{Mistrial}}
Declaración formal de nulidad del juicio. El prefijo \textit{mis-} indica algo erróneo: el juicio no sirve y hay que iniciarlo de nuevo. Un \textit{hung jury} puede conducir a un \textit{mistrial}.
\end{definition}

Cuando se declara el \textit{mistrial} no hubo sentencia todavía; se convoca un nuevo jurado y se vuelve a la etapa de alegatos y siguientes. La persona ya está procesada, pero todo se pone nuevamente en marcha.

\begin{table}[H]
\centering
\small
\caption{\textit{Hung jury} frente a \textit{mistrial}}
\begin{tabularx}{\textwidth}{L L}
\toprule
\textbf{\textit{Hung jury}} & \textbf{\textit{Mistrial}} \\
\midrule
Jurado bloqueado: no logra unanimidad ni para condenar ni para absolver. & Declaración formal de nulidad del juicio. \\
\addlinespace
Es la situación que se anuncia al juez: ``no pudimos llegar a un veredicto''. & Consecuencia posible: se convoca un nuevo jurado y se retoma desde la etapa probatoria. \\
\addlinespace
No importa cuántos discrepen (puede ser 11 a 1). & No hubo sentencia todavía; el juicio ``no sirve'' y se reinicia. \\
\bottomrule
\end{tabularx}
\end{table}

\begin{example}{Caso reciente comentado en clase}
Una estudiante recordó el caso de una mujer que mató a sus tres hijos, intentó suicidarse y quedó paralizada. Se discutía si estaba en plena posesión de sus facultades, pues tenía depresión posparto. Según una nota leída en clase, el jurado de 12 personas deliberó durante siete días (alrededor de 28 horas) sin llegar a un consenso. La votación estaba 11 contra uno: 11 integrantes apoyaban la absolución por motivos de salud mental, pero un solo jurado se negó a cambiar su voto. Al no existir unanimidad, la jueza declaró la nulidad del juicio. % TODO: las cifras difieren entre el relato oral (siete jurados con dudas) y la nota leída (11 contra uno)
\end{example}

\subsection{Costo, contaminación del jurado y renuncia al jurado}

El docente destacó que se trata de un procedimiento muy costoso y engorroso. En una causa con repercusión mediática, además, surge el problema de dónde obtener un nuevo jurado que no esté contaminado. Normalmente los miembros del jurado son de la localidad donde está el juzgado, pero en un caso de mucha repercusión pueden convocarse jurados de otro lugar o de localidades vecinas para que no estén tan empapados del caso.

Una estudiante señaló que los nuevos jurados ya habrán consumido noticias sobre el caso y podrían tener preconceptos, por lo que una alternativa para la acusada sería renunciar al derecho al jurado y que resuelva un juez. El docente confirmó que puede renunciarse a ese derecho. Agregó que, aun sin cobertura mediática, es muy difícil mantener al jurado incontaminado, porque las personas comentan lo que las impacta; la única forma es aislarlo, pero es difícil hacerlo por mucho tiempo en la actualidad, ya que las redes sociales hacen que todo pueda viralizarse.

Como comparación, en el sistema argentino también se puede anular una causa y empezar de nuevo, pero, al ser un proceso escrito o mayormente escrito, el docente lo ve como algo más leve. Mencionó como ejemplo un caso mediático local en el que el juicio estaba viciado y hubo que reiniciarlo. % TODO: el caso local mencionado en clase no está identificado en las notas originales

\section{Cuadro Cornell: el juicio y el jurado}

\begin{longtable}{@{}>{\raggedright\arraybackslash}p{0.28\textwidth}>{\raggedright\arraybackslash}p{0.66\textwidth}@{}}
\toprule
\textbf{Preguntas / palabras clave} & \textbf{Notas / respuestas} \\
\midrule
\endfirsthead
\toprule
\textbf{Preguntas / palabras clave} & \textbf{Notas / respuestas} \\
\midrule
\endhead
\bottomrule
\endlastfoot

\textbf{¿Cómo se desarrolla el \textit{trial}?} &
Igual que en el proceso civil: alegatos de apertura, producción de prueba y alegatos de cierre (15 minutos por parte). Comienza la fiscalía; en el cierre, la defensa habla primero. Se convocan más de 12 jurados y suplentes. \\
\addlinespace

\textbf{¿Por qué puede haber tres jurados?} &
\textit{Grand jury} (acusar), \textit{petit jury} (valorar la prueba y decidir la culpabilidad) y \textit{sentencing jury} (determinar la pena en delitos gravísimos). \\
\addlinespace

\textbf{¿Qué hace el juez y qué hace el jurado?} &
El juez controla las reglas de juego y lo jurídico (admisibilidad, instrucciones); el jurado evalúa lo fáctico. Las instrucciones buscan neutralizar influencias, aunque ``borrar'' lo escuchado es muy difícil. \\
\addlinespace

\textbf{¿Qué estándar rige y cuántos votos hacen falta?} &
Penal: \textit{beyond a reasonable doubt} y unanimidad de 12. Civil: preponderancia de la prueba y no necesariamente unanimidad (por ejemplo, 4 de 6). \\
\addlinespace

\textbf{¿Qué es un \textit{hung jury} y un \textit{mistrial}?} &
\textit{Hung jury}: jurado bloqueado, sin veredicto. \textit{Mistrial}: declaración formal de nulidad; se convoca un nuevo jurado y se reinicia el \textit{trial}. Es costoso y plantea el riesgo de jurados contaminados por los medios. \\

\midrule
\multicolumn{2}{@{}p{\dimexpr0.94\textwidth+2\tabcolsep\relax}@{}}{%
\textbf{Resumen:} El \textit{trial} reproduce la estructura civil (apertura, prueba, cierre), pero con un estándar mucho más exigente (más allá de toda duda razonable) y un jurado de 12 que debe decidir por unanimidad. El juez controla lo jurídico y el jurado valora los hechos. Si el jurado no logra un veredicto (\textit{hung jury}), se declara el \textit{mistrial} y se reinicia el juicio con un nuevo jurado.} \\
\end{longtable}

% ============================================================
\chapter[Sentencia y determinación de la pena]{Sentencia y determinación de la pena}
% ============================================================

\section{Penas posibles (\textit{sentencing})}

Suponiendo que no se declare el \textit{mistrial}, el jurado emite un veredicto que será \textit{guilty} o \textit{not guilty}. Lo que sigue es la \textit{sentencing stage} o \textit{sentencing phase}.

\begin{important}
\textit{Sentence} no significa ``sentencia'': es un falso amigo. \textit{Sentencing} es la determinación de la pena o sanción. Según el texto, normalmente la sanción se impone \textit{immediately} en las infracciones, cuando son delitos menores: enseguida después de que el jurado dice ``sí''.
\end{important}

Las posibles penas, ya vistas en otro texto, son las siguientes:

\begin{itemize}
    \item \textit{Fines}: multas.
    \item Prisión por un período breve.
    \item Prisión por un período más prolongado.
    \item \textit{Suspended sentence}: la pena existe, pero su ejecución está en suspenso y solo se ejecutará si no se cumplen las condiciones impuestas (\textit{to conditions}).
    \item \textit{Payment of restitution to the victim}: pago de algún tipo de compensación o retribución a la víctima.
    \item \textit{Community service}.
    \item \textit{Alcohol rehabilitation}.
    \item \textit{Probation}.
\end{itemize}

\subsection{\textit{Probation}: Estados Unidos frente a Argentina}

En Estados Unidos la \textit{probation} se incluye como una de las posibles sanciones o condenas, mientras que para nosotros no es una pena: lo que aquí se llama (impropiamente) \textit{probation} está mucho más adelante en el proceso, pues se suspende el juicio a prueba y el proceso no sigue avanzando. En Estados Unidos, en cambio, la persona ya atravesó todo el proceso, ya fue declarada culpable, y una de las posibilidades es que se le imponga una \textit{probation}.

\begin{table}[H]
\centering
\small
\caption{\textit{Probation}: Estados Unidos frente a Argentina}
\begin{tabularx}{\textwidth}{>{\raggedright\arraybackslash}p{0.24\textwidth} L L}
\toprule
\textbf{Aspecto} & \textbf{Estados Unidos} & \textbf{Argentina} \\
\midrule
\textbf{Naturaleza} & Una de las posibles penas o condenas. & No es una pena: es la suspensión del juicio a prueba. \\
\addlinespace
\textbf{Momento del proceso} & Después del veredicto de culpabilidad. & Mucho antes: se suspende el juicio y el proceso no avanza. \\
\addlinespace
\textbf{Situación del acusado} & Ya fue declarado culpable. & Aún no hay condena. \\
\bottomrule
\end{tabularx}
\end{table}

\section{El \textit{sentencing jury} y el ``mini juicio''}

En algunas causas por delitos graves puede convocarse un nuevo jurado, el \textit{sentencing jury}: ya hubo un \textit{grand jury}, luego el \textit{petit jury}, y ahora el \textit{sentencing jury}. El texto no indica cuántas personas lo integran. % TODO: contenido incompleto o ambiguo en las notas originales (cantidad de integrantes del sentencing jury)

Este jurado ya no vuelve a analizar si la persona es culpable, porque eso ya fue resuelto; en función de la prueba, debe determinar la pena que corresponde aplicar. Según el texto, en otras jurisdicciones esta decisión la toma el juez para los delitos graves, solución que al docente le parece la más lógica desde su perspectiva. % TODO: contenido incompleto o ambiguo en las notas originales (alcance de lo que dice el texto sobre otras jurisdicciones)

\begin{definition}{\textit{Sentencing hearing} (``mini juicio'')}
Audiencia destinada no a evaluar la culpabilidad sino a evaluar qué pena corresponde. La fiscalía presenta evidencia de \textit{aggravating factors} (circunstancias agravantes) y la defensa presenta los \textit{mitigating factors} (circunstancias atenuantes).
\end{definition}

\begin{note}
Según el docente, este jurado no necesita estar asesorado por un abogado. Su tarea consiste en determinar si, por las circunstancias planteadas (por ejemplo, que el acusado atravesaba una depresión o era víctima de violencia, y todos los agravantes y atenuantes presentados), corresponde la pena de muerte o la prisión perpetua.
\end{note}

Una estudiante informó, a partir de una búsqueda en internet, que en casos de pena de muerte o prisión perpetua el mismo jurado de 12 personas suele deliberar en una segunda fase. El docente respondió que habría que ver las jurisdicciones: el texto dice \textit{with a completely new jury}, y para verificarlo habría que consultar el \textit{U.S. Code}, que es federal y contiene recomendaciones, y luego algunos estados al azar, ya que cada estado tiene su propio código procesal. Reconoció tener dudas, aunque lo tenía por un nuevo jurado. % TODO: contenido ambiguo en las notas originales (mismo jurado o jurado nuevo)

\begin{formula}{Interpretación de la Corte Suprema (según el texto leído en clase)}
La Constitución protege \textit{``the right to a jury sentencing trial for all defendants facing the death penalty''}.

Traducción: la Constitución protege el derecho a un juicio de sentencia por jurado para todo acusado que enfrente la pena de muerte. Según lo expuesto, cuando la pena posible es la pena de muerte, la Constitución, tal como la interpretó la Corte Suprema, contempla este mini juicio para determinar si corresponde imponer esa pena.
\end{formula}

En definitiva, la pena la determina directamente el juez o este jurado, pero lo que se decide no es la sentencia, sino la determinación de la pena.

\section{La \textit{allocution}}

Antes de que el juez anuncie la pena (la haya resuelto el jurado o el propio juez), el texto indica que se da al acusado la oportunidad de la \textit{allocution}: el condenado puede o no hablar, y lo hace directamente ante el juez, sin su abogado. Es una última oportunidad: antes de que el juez anuncie la pena, se le pregunta qué tiene para decir.

\begin{definition}{\textit{Allocution}}
Última oportunidad que se da al condenado de dirigirse directamente al juez, sin intermediación de su abogado, antes de que se anuncie la pena. Puede optar por no decir nada.
\end{definition}

A lo largo de todo el proceso el acusado nunca se dirige directamente al juez ni puede hablar a los miembros del jurado: siempre actúa a través de los abogados, y el juez lo impide, pues un acusado de un delito grave que se dirigiera al jurado sería muy intimidante. Por eso la \textit{allocution} se ubica poco antes de la sentencia.

Según el texto, sirve para demostrar arrepentimiento (\textit{show remorse}) y para ofrecer las motivaciones detrás de los actos criminales (\textit{offer the motivations behind the criminal acts}). Nada de lo que diga cambiará el veredicto, pero puede aportar algo que atenúe la pena. Según el texto, puede además mostrar hechos que nadie conocía y que no estaban expuestos durante el proceso, de modo que la pena podría variar por lo dicho. No obstante, el docente considera que es más simbólica que otra cosa: responde a cuestiones morales, religiosas o relativas a la víctima si hay arrepentimiento, y puede ocurrir que lo que se diga empeore la situación y afecte aún más a los familiares de las víctimas. Para él es una formalidad, aunque es un derecho que existe.

Respecto del sistema argentino, el docente cree que esta figura no existe: el acusado habla antes de que se cierre todo, como último en hablar, pero todavía no está condenado. Durante el debate el acusado tiene derecho a declarar, pero la \textit{allocution} apunta a otra cosa: es el único momento del proceso en que puede obviar a su abogado y hablar directamente al juez, cuando todo lo demás debe encauzarlo a través del abogado, como estrategia procesal.

\begin{important}
La \textit{allocution} es el único momento del proceso en que el acusado puede hablarle directamente al juez, sin intermediación de su abogado, antes de que se anuncie la pena. Es un derecho consagrado, aunque el docente lo considera más simbólico que decisivo.
\end{important}

\section{Cierre del bloque procesal}

Salvo toda la parte de las distintas audiencias (en particular, el \textit{plea bargaining}, visto en la clase anterior), el resto del proceso penal estadounidense es muy parecido al civil, sin demasiadas diferencias. Es importante comprender la diferencia del estándar probatorio y de la cantidad de miembros del jurado; la forma en que se lleva adelante el proceso es bastante similar.

\section{Cuadro Cornell: sentencia y determinación de la pena}

\begin{longtable}{@{}>{\raggedright\arraybackslash}p{0.28\textwidth}>{\raggedright\arraybackslash}p{0.66\textwidth}@{}}
\toprule
\textbf{Preguntas / palabras clave} & \textbf{Notas / respuestas} \\
\midrule
\endfirsthead
\toprule
\textbf{Preguntas / palabras clave} & \textbf{Notas / respuestas} \\
\midrule
\endhead
\bottomrule
\endlastfoot

\textbf{¿Qué penas existen?} &
Multas, prisión breve o prolongada, \textit{suspended sentence}, restitución a la víctima, servicio comunitario, rehabilitación por alcohol y \textit{probation} (que en Estados Unidos es una pena, no una suspensión del juicio). \\
\addlinespace

\textbf{¿Qué es el \textit{sentencing jury}?} &
Un ``mini juicio'' para fijar la pena: la fiscalía presenta agravantes y la defensa atenuantes. La Constitución, según la Corte Suprema, protege este derecho en casos de pena de muerte. Hay dudas sobre si es el mismo jurado u otro. \\
\addlinespace

\textbf{¿Qué es la \textit{allocution}?} &
Única oportunidad en que el condenado habla directamente al juez, sin su abogado, antes de anunciarse la pena: puede mostrar arrepentimiento o aportar motivaciones. Es un derecho, más simbólico que decisivo según el docente. \\

\midrule
\multicolumn{2}{@{}p{\dimexpr0.94\textwidth+2\tabcolsep\relax}@{}}{%
\textbf{Resumen:} Tras el veredicto se determina la pena (\textit{sentencing}, no ``sentencia''), que puede ser multa, prisión, \textit{suspended sentence}, restitución, servicio comunitario, rehabilitación o \textit{probation}. En delitos gravísimos puede intervenir un tercer jurado, mediante un ``mini juicio'' de agravantes y atenuantes. Antes de anunciarse la pena, el condenado conserva la \textit{allocution} como última palabra directa ante el juez.} \\
\end{longtable}

% ============================================================
\chapter[Garantías constitucionales]{Garantías constitucionales: lectura del texto}
% ============================================================

\section{Presentación del texto y debido proceso}

El texto de lectura comienza hablando de las enmiendas de la Constitución (\textit{the Amendments of the Constitution}). Presenta términos resaltados y fallos, y menciona, entre otros, \textit{public trial}, \textit{right to a speedy trial} y \textit{witness against himself}, así como el \textit{Bill of Rights}, \textit{Miranda v. Arizona}, \textit{search warrant}, \textit{Federal Rules of Criminal Procedure} y la Constitución. En lugar de explicar las etapas del proceso, el texto explica sus fases, desglosando las enmiendas, tomando sus partes, ampliándolas y señalando en qué contexto se comprende cada instituto. Hay abundante jurisprudencia.

\begin{definition}{\textit{Search warrant} / \textit{arrest warrant}}
Orden judicial de allanamiento o de detención, para que nadie sea detenido ni allanado sin una debida orden judicial.
\end{definition}

Las garantías mencionadas (por ejemplo, que no se puede exigir a una persona que declare contra sí misma, o el \textit{indictment}) existen para que haya un debido proceso. Se trata de todas las garantías del debido proceso, entendido como un derecho amplio que abarca muchas cosas.

\begin{definition}{Debido proceso (\textit{due process of law})}
Derecho amplio que integra numerosas garantías. No es una cosa sola: es un derecho mucho más grande del que forman parte todos los derechos identificados en clase.
\end{definition}

\begin{table}[H]
\centering
\small
\caption{Garantías que integran el debido proceso (identificadas en clase)}
\begin{tabularx}{\textwidth}{>{\raggedright\arraybackslash}p{0.28\textwidth} L}
\toprule
\textbf{Garantía} & \textbf{Explicación} \\
\midrule
\textbf{Juez natural y ley anterior} & El tribunal y las reglas deben existir antes del hecho. \\
\addlinespace
\textbf{Derecho de defensa} & Ser escuchado, presentar testigos y contar con abogado. \\
\addlinespace
\textbf{Defensa técnica y eficiente} & No basta con tener abogado: debe actuar con diligencia (plazos, pedidos, testigos, pericias). \\
\addlinespace
\textbf{Juez imparcial} & Imparcialidad de quien juzga. \\
\addlinespace
\textbf{Juicio por jurado} & Derecho a optar por jurado en los casos en que corresponde. \\
\addlinespace
\textbf{Garantías policiales y de prueba} & Orden judicial (\textit{warrant}), advertencias, cadena de custodia. \\
\bottomrule
\end{tabularx}
\end{table}

\subsection{Defensa técnica y eficiente}

No solo se requiere que una persona sea representada por un abogado y tenga su debida defensa, sino que el abogado que la representa haga las cosas bien. En algunas causas se planteó que el acusado no había tenido una defensa técnica y además eficiente. Esto no se debe a que el resultado no haya sido favorable o a que no se haya logrado la absolución, sino a que el abogado no había hecho lo suficiente: se le pasaban los plazos, no contestaba los pedidos, no buscaba testigos, no confrontaba. En apelación, el acusado sostuvo que no tuvo una defensa sólida.

\begin{example}{Novela comentada en clase}
El docente recordó una novela sobre un hombre condenado a pena de muerte por un delito de violación seguido de homicidio, que esperó mucho tiempo la ejecución de la pena. Asociaciones y ONG contrarias a la pena de muerte, que utilizan por ejemplo el ADN y ayudan a revisar causas que de otro modo quedarían archivadas, reactivan la causa y determinan varios vicios procesales: por ejemplo, que al jurado no se le habían dado instrucciones. Además, el abogado defensor nunca pidió una pericia psicológica o psiquiátrica, y la persona resultó ser incapaz: un adulto con la comprensión de un niño de siete años en un juicio de adultos, con retraso madurativo, que sonreía mientras se describían los hechos, lo que llevó al jurado a considerarlo perverso. Se concluyó que no había tenido una defensa sólida, válida y eficiente, porque ni siquiera se pidió una pericia que era elemental dadas las circunstancias.
\end{example}

\section{La Quinta Enmienda}

La redacción de la Quinta Enmienda es antigua y excesivamente formal; consiste en una oración larguísima, algo raro en inglés. Utiliza \textit{shall}, verbo modal que indica la idea de una obligación impuesta por una autoridad.

\begin{formula}{Quinta Enmienda (texto de referencia)}
\textit{No person shall be held to answer for a capital, or otherwise infamous crime, unless on a presentment or indictment of a Grand Jury, except in cases arising in the land or naval forces, or in the Militia, when in actual service in time of War or public danger; nor shall any person be subject for the same offence to be twice put in jeopardy of life or limb; nor shall be compelled in any criminal case to be a witness against himself, nor be deprived of life, liberty, or property, without due process of law; nor shall private property be taken for public use, without just compensation.}

\medskip
\textbf{Traducción:} Nadie será obligado a responder por un delito capital o infame sin una acusación formal (\textit{presentment} o \textit{indictment}) de un gran jurado, salvo en casos que surjan en las fuerzas terrestres o navales, o en la milicia, estando en servicio activo en tiempo de guerra o peligro público; ni persona alguna será sometida dos veces, por el mismo delito, a peligro de perder la vida o un miembro; ni será obligada, en ningún proceso penal, a declarar contra sí misma; ni será privada de la vida, la libertad o la propiedad sin el debido proceso legal; ni se tomará propiedad privada para uso público sin justa compensación.

\medskip
\textit{El texto original de la enmienda no fue leído de corrido por el docente; se reproduce como referencia, y a continuación se expone su explicación.}
\end{formula}

\subsection{Derechos consagrados}

\begin{enumerate}
    \item \textbf{Acusación formal previa.} No se puede juzgar a ninguna persona por un delito grave o infame (aberrante) salvo que exista un \textit{indictment} o \textit{presentment} (también se usa \textit{information}): una acusación formal de un juez o de un gran jurado. Sin acusación formal no se puede someter a una persona a procesamiento.
    \item \textbf{Excepción.} Causas en tiempos de guerra o relacionadas con las fuerzas terrestres o navales, o con la milicia. Se entiende que, al redactarse las enmiendas, no existían ejércitos profesionales sino la organización de fuerzas revolucionarias o fuerzas que defendían al país. La expresión \textit{when in actual service in time of war or public danger} alude a épocas de guerra o de amenaza pública, aunque no haya una amenaza concreta pero estén por invadir; en esos casos se puede actuar sin cumplir todas estas garantías. Un estudiante mencionó como ejemplo el de un traidor a la patria: un soldado que ayuda a un ejército extranjero. La regla es que no puede haber procesamiento sin acusación formal, salvo en momentos de guerra o peligro público, o cuando se trate de las fuerzas armadas.
    \item \textbf{\textit{Double jeopardy}.} La oración continúa con un punto y coma y \textit{nor shall} (``ni tampoco''). Nadie puede ser sometido dos veces, por el mismo delito, a \textit{jeopardy of life or limb}. \textit{Limb} significa brazo, pierna, extremidades; la expresión viene de otra época, en la que quizá se aplicaban castigos corporales, y hoy resulta obsoleta. La garantía impide procesar a una persona dos veces por la misma circunstancia o delito.
    \item \textbf{No autoincriminación.} \textit{Compelled} da la idea de obligar, de forzar a una persona a declarar contra sí misma en una causa penal.
    \item \textbf{Debido proceso.} Nadie puede ser privado de la vida (ejecutar la pena de muerte), de la libertad (asignar pena de prisión) ni de la \textit{property} sin un debido proceso. \textit{Property} se entiende como término amplio: bienes en general, no solo inmuebles, a diferencia de nuestro uso de ``propiedad''.
    \item \textbf{Justa compensación.} La propiedad privada no puede tomarse para uso público sin compensación: no se la puede expropiar sin una debida compensación o retribución.
\end{enumerate}

\begin{example}{``I plead the Fifth''}
En las películas, cuando se hace declarar al acusado, aparece la frase \textit{I plead the Fifth}, que a veces se traduce como ``me atengo'' o ``me adhiero a la Quinta''. Significa que no se está obligado a declarar: no pueden obligarlo si siente que lo que va a decir lo incriminaría. Es el derecho a guardar silencio; no se dice más que eso y es suficiente.
\end{example}

Una estudiante preguntó si podría invocarse esta enmienda por un secreto profesional (por ejemplo, de un médico). El docente respondió que en ese caso se diría que se está resguardado por el secreto profesional y no se invocaría esta enmienda, que habla concretamente de no declarar contra uno mismo. Admitió, sin embargo, que si el hecho puede generar un grave perjuicio a quien declara, podría ser distinto. La cláusula apunta más directamente al acusado que no quiere colaborar con la fiscalía y darle argumentos.

\section{La Sexta Enmienda}

\begin{formula}{Sexta Enmienda (texto de referencia)}
\textit{In all criminal prosecutions, the accused shall enjoy the right to a speedy and public trial, by an impartial jury of the State and district wherein the crime shall have been committed, which district shall have been previously ascertained by law, and to be informed of the nature and cause of the accusation; to be confronted with the witnesses against him; to have compulsory process for obtaining witnesses in his favor, and to have the Assistance of Counsel for his defence.}

\medskip
\textbf{Traducción:} En todos los procesos penales, el acusado gozará del derecho a un juicio rápido y público, ante un jurado imparcial del Estado y distrito en que se haya cometido el delito, distrito que deberá haber sido determinado previamente por ley; a ser informado de la naturaleza y causa de la acusación; a ser confrontado con los testigos en su contra; a contar con un procedimiento compulsivo para obtener testigos a su favor; y a contar con la asistencia de un abogado para su defensa.

\medskip
\textit{Texto de referencia; en clase se comentó cada uno de estos elementos.}
\end{formula}

Derechos que contiene:

\begin{itemize}
    \item \textbf{Juicio rápido y público} (\textit{speedy and public trial}).
    \item \textbf{Jurado imparcial} (\textit{impartial jury}): el texto no dice \textit{judge}, dice \textit{jury}. El jurado debe ser del estado y distrito donde se hubiera cometido el delito. El distrito ya tiene que existir desde antes: no puede armarse un distrito después de cometido el hecho para encausar una causa. Se lo puede pensar desde la idea de juez natural y ley anterior: todo debe existir antes de la comisión del delito.
    \item \textbf{Información de la acusación.} Se informa al acusado de la naturaleza y la causa de la acusación; ese derecho se garantiza en la audiencia de \textit{arraignment}, en que formalmente se le comunica de qué se lo acusa.
    \item \textbf{Confrontación.} Derecho a enfrentarse a los testigos que declaran en su contra, tener la oportunidad de confrontarlos y saber quiénes declararán.
    \item \textbf{Prueba compulsiva de testigos propios} (\textit{compulsory process}): que la justicia colabore para que quien fue convocado como testigo, o para producir un peritaje, se presente a declarar, mediante la intervención de la fuerza pública o de la autoridad judicial. El acusado puede pedir al juzgado ayuda porque necesita a determinada persona como testigo y esta no se presenta, y se le llevará la intimación.
    \item \textbf{Asistencia de abogado} para su defensa.
\end{itemize}

\section{Reflexiones sobre las Enmiendas Quinta y Sexta}

Estas dos enmiendas son las más relevantes. El derecho al juicio por jurado está consagrado en la Constitución, igual que en el caso argentino. Sin embargo, en Estados Unidos es la regla, mientras que en la Argentina no se hizo efectivo hasta hace muy poco tiempo y no en todas las causas. Es una idea muy valorada del sistema anglosajón: el \textit{jury of peers}, un jurado de pares. La sociedad lo ve como algo sumamente positivo: no es una autoridad que viene desde arriba, alejada de la situación, quien resuelve el destino de una persona, sino un grupo de personas como cualquiera, con todo lo bueno y lo malo.

Además, deben recordarse estas ideas: el distrito donde se lleva adelante la causa ya tiene que estar determinado antes; el juicio debe ser rápido y público; y debe existir una acusación formal, \textit{unless} (salvo que haya) \textit{presentment} o \textit{indictment}.

\begin{definition}{\textit{Indictment}, \textit{presentment} e \textit{information}}
\textit{Indictment} es lo que emite el gran jurado cuando decide acusar. \textit{Presentment} es lo mismo. Cuando la acusación se hace directamente, sin un gran jurado, se denomina \textit{information}: es igualmente un escrito formal.
\end{definition}

\section{Cuadro Cornell: garantías constitucionales}

\begin{longtable}{@{}>{\raggedright\arraybackslash}p{0.28\textwidth}>{\raggedright\arraybackslash}p{0.66\textwidth}@{}}
\toprule
\textbf{Preguntas / palabras clave} & \textbf{Notas / respuestas} \\
\midrule
\endfirsthead
\toprule
\textbf{Preguntas / palabras clave} & \textbf{Notas / respuestas} \\
\midrule
\endhead
\bottomrule
\endlastfoot

\textbf{¿Qué integra el debido proceso?} &
Juez natural, ley anterior, derecho de defensa, ser oído, juez imparcial, juicio por jurado y defensa técnica eficiente (como en la novela del condenado sin pericia psiquiátrica). \\
\addlinespace

\textbf{¿Qué protege la Quinta Enmienda?} &
Acusación formal previa (gran jurado), prohibición de \textit{double jeopardy}, derecho a no declarar contra sí mismo (``\textit{I plead the Fifth}''), debido proceso y justa compensación por expropiación. \\
\addlinespace

\textbf{¿Qué garantiza la Sexta Enmienda?} &
Juicio rápido y público, jurado imparcial del distrito preexistente, información de la acusación, confrontación de testigos, compulsión de testigos propios y asistencia de abogado. \\
\addlinespace

\textbf{¿Qué diferencia hay entre \textit{indictment}, \textit{presentment} e \textit{information}?} &
\textit{Indictment}: acusación formal que emite el gran jurado. \textit{Presentment}: lo mismo. \textit{Information}: acusación formal escrita que se hace directamente, sin gran jurado. \\

\midrule
\multicolumn{2}{@{}p{\dimexpr0.94\textwidth+2\tabcolsep\relax}@{}}{%
\textbf{Resumen:} El debido proceso es un derecho amplio, anclado en las Enmiendas Quinta y Sexta, que reúne garantías como la acusación formal previa, la prohibición del doble juzgamiento, el derecho a no declarar contra sí mismo, el juicio rápido y público ante un jurado imparcial de pares, la confrontación de testigos y la asistencia de un abogado, que además debe ejercer una defensa eficiente.} \\
\end{longtable}

% ============================================================
\chapter[Fuentes normativas y procedimiento policial]{Fuentes normativas y procedimiento policial}
% ============================================================

\section{El \textit{U.S. Code} y los códigos estatales}

El texto menciona el \textit{U.S. Code} y señala que los distintos distritos tienen sus propios códigos de derecho penal: municipalidades, estados y gobierno federal. Las distintas subdivisiones del país tienen cada una sus propios códigos de fondo y de forma, a diferencia de lo que sucede en la Argentina, donde rige un único Código Penal en todo el país.

El docente advirtió sobre una falacia: muchos creen que no hay un código a nivel federal en Estados Unidos, y sí lo hay: el \textit{U.S. Code}. Lo describió como una ``bolsa de gatos'', porque en él se mezclan todas las áreas. El título 18 (según lo dicho en clase) se refiere a todo lo penal: la tipificación de los delitos y lo relativo al proceso. Otro título, el 11 (según lo recordado por el docente), trata de concursos y quiebras, lo que lleva a lo comercial; otro se refiere a lo marítimo, el derecho de almirantazgo y la Marina. % TODO: contenido incompleto o ambiguo en las notas originales (numeración de los títulos 18 y 11 dicha con dudas)

\begin{formula}{\textit{U.S. Code}: títulos mencionados en clase (orientativo, no obligatorio)}
Título 2: Congreso. Título 3: Presidente. Título 4: Bandera. Título 5: Organización del gobierno y empleados. Título 6: Seguridad nacional. Título 7: Agricultura. Título 8: Extranjería. Título 9: Arbitraje. Título 10: Fuerzas armadas. Título 11: Quiebras. Título 12: Bancos.

\medskip
El docente señala que cada título (\textit{title}) equivale, para nosotros, a un libro o capítulo, y que cada uno aborda un tema en particular. Los primeros títulos tienen que ver con la organización del país y del Estado; después se mezclan materias (por ejemplo, comercial y quiebras en el título 11).
\end{formula}

El \textit{U.S. Code} es de nivel federal, pero constituye una recomendación: los estados no están obligados a aplicarlo, a diferencia de lo que ocurre en la Argentina con el Código Penal. Es orientativo, a los fines de que cada estado o subdivisión adopte su propio código, que en general lo siguen. Sin embargo, existen diferencias enormes, por ejemplo con la pena de muerte: hay estados donde está contemplada y otros donde no, de modo que no hay un parámetro uniforme.

\begin{table}[H]
\centering
\small
\caption{Sistema normativo penal: Argentina frente a Estados Unidos}
\begin{tabularx}{\textwidth}{>{\raggedright\arraybackslash}p{0.22\textwidth} L L}
\toprule
\textbf{Aspecto} & \textbf{Argentina} & \textbf{Estados Unidos} \\
\midrule
\textbf{Código de fondo} & Un Código Penal único para todo el país. & Cada estado tiene el suyo; el \textit{U.S. Code} es federal y orientativo. \\
\addlinespace
\textbf{Código procesal} & Hay uniformidad; cada provincia puede tener el suyo, pero sin apartarse radicalmente. & \textit{Federal Rules of Criminal Procedure} (orientativas) y reglas propias de cada estado. \\
\addlinespace
\textbf{Pena de muerte} & No aplicable. & Algunos estados la contemplan y otros no. \\
\addlinespace
\textbf{Conductas que son delito} & Las mismas en todo el país. & Pueden variar: lo que es delito en un lugar puede no serlo en otro. \\
\bottomrule
\end{tabularx}
\end{table}

\subsection{Normas de procedimiento}

Lo mismo ocurre con las normas de procedimiento. Las \textit{Federal Rules of Criminal Procedure} equivalen a nuestro código procesal: establecen en qué casos hay que excluir la prueba, si es válido el testimonio de oídas, qué sucede cuando un testigo no quiere comparecer, y otras pautas del proceso a nivel federal; después, cada estado tiene las propias.

Ante la pregunta de si hay algún límite, el docente respondió que debe haberlo, aunque el sistema es muy distinto del nuestro, más unificado: en la Argentina cada provincia tiene su propio código procesal, pero no se aparta radicalmente, porque hay un presupuesto mínimo del procesal nacional y, por delegación, cada provincia puede plantear su propio procedimiento, con bastante uniformidad. Para el docente, más que la cuestión de forma, lo más extraño para nuestra mirada es el código de fondo: que en un lugar estén contempladas penas que en otro no, o que haya conductas que son delito en un lugar y en otro no.

\begin{important}
Para investigar, el \textit{U.S. Code} sirve como orientación (``a ver qué pasa en Estados Unidos con tal figura''), pero hay que tener siempre en cuenta que no es obligatorio y que no rige para todo el país: si se quiere buscar algo en particular de un estado, hay que ir a las normas de ese estado.
\end{important}

\section{\textit{Federal Rules} y el ``piso mínimo''}

Las normas federales siguientes son en su mayoría orientativas, y si los estados deciden apartarse, se apartan. Las \textit{Federal Rules} recogen los derechos consagrados en el \textit{Bill of Rights}, los que se vieron en la Constitución, y constituyen un piso mínimo, porque la Constitución incluye estas garantías.

\begin{formula}{Garantías protegidas por la Constitución y las \textit{Federal Rules} (según el texto)}
\textit{Equal protection under the laws; the right to have legal counsel; the right to confront witnesses; the right to a trial; the right not to testify against oneself.}

\textit{State constitutions and procedural rules may increase protection (``increase protection'') above what the Federal Constitution provides.}

\medskip
\textbf{Traducción:} protección igualitaria ante la ley; derecho a contar con asistencia letrada; derecho a confrontar testigos; derecho a un juicio; derecho a no declarar contra uno mismo. Las constituciones y reglas procesales estatales pueden aumentar la protección por encima de lo que dice la Constitución federal.

\medskip
\textbf{Aclaración del docente:} el texto es muy repetitivo y vuelve a mencionar todo: por encima, más garantías, más protecciones; por debajo, no.
\end{formula}

Al investigar algo puede usarse el \textit{U.S. Code} como orientación, pero no es obligatorio. El piso mínimo está dado por la Constitución: los estados pueden agregar garantías, pero no quitarlas.

\section{Procedimiento policial, advertencias Miranda y debido proceso sustantivo}

El texto trata luego lo que sucede a nivel policial, antes de llegar a la instancia judicial. Menciona los \textit{Miranda Rights}, referidos al procedimiento de detención de una persona y a las garantías que tiene durante el arresto, antes de que llegue al juzgado. Todo lo que ocurre en sede policial forma parte de impartir justicia, y se busca velar porque ese procedimiento cumpla con las garantías individuales.

\begin{formula}{\textit{Miranda v. Arizona} y las pautas policiales (texto leído en clase)}
\textit{``Provide guidance and rules that law enforcement must follow.''}

\textit{``Failure to follow such procedure may [result in the exclusion of evidence or the release] of arrested suspect.''}

\medskip
\textbf{Traducción:} brindan pautas y reglas que las fuerzas de seguridad deben seguir. No seguir tal procedimiento puede acarrear [la exclusión de pruebas o la liberación] del sospechoso detenido.

\medskip
\textit{La grabación es parcialmente inaudible en este pasaje; la lectura del docente se completa según su explicación posterior.} % TODO: contenido incompleto o ambiguo en las notas originales
\end{formula}

El texto habla de las pautas que la policía debe seguir para que el procedimiento se lleve a cabo conforme a derecho. Si no se siguen las pautas procesales como corresponde, puede ocurrir que una prueba contundente y determinante quede excluida o que incluso la persona quede en libertad.

El texto menciona también el \textit{substantive due process} (debido proceso sustantivo). Un acusado penal debe conocer sus derechos de antemano: si va a declarar algo, debe saber que el Estado puede usar luego esas declaraciones; lo que diga en la instancia de la detención puede ser utilizado después por la fiscalía. Y debe haber entendido: no alcanza con que se le diga como un trámite, como si se lo enunciara, sino que debe haber entendido y comprendido. % TODO: contenido incompleto o ambiguo en las notas originales (la grabación se interrumpe en este punto)

\begin{important}
No alcanza con enunciar los derechos del detenido como un mero trámite: la persona debe haberlos entendido y comprendido, porque lo que declare puede ser usado luego por la fiscalía. Si la policía no sigue el procedimiento, la prueba obtenida puede quedar excluida e incluso el sospechoso puede recuperar la libertad.
\end{important}

\section{Cuadro Cornell: fuentes normativas y procedimiento policial}

\begin{longtable}{@{}>{\raggedright\arraybackslash}p{0.28\textwidth}>{\raggedright\arraybackslash}p{0.66\textwidth}@{}}
\toprule
\textbf{Preguntas / palabras clave} & \textbf{Notas / respuestas} \\
\midrule
\endfirsthead
\toprule
\textbf{Preguntas / palabras clave} & \textbf{Notas / respuestas} \\
\midrule
\endhead
\bottomrule
\endlastfoot

\textbf{¿Rige el \textit{U.S. Code} en todo el país?} &
No: es federal y orientativo. Cada estado tiene su propio código de fondo y de forma; difieren penas y conductas (por ejemplo, la pena de muerte). \\
\addlinespace

\textbf{¿Qué son las \textit{Federal Rules} y qué es el ``piso mínimo''?} &
Las \textit{Federal Rules of Criminal Procedure} equivalen a nuestro código procesal a nivel federal y recogen los derechos del \textit{Bill of Rights}. El piso mínimo lo da la Constitución: los estados pueden ampliar la protección, pero no reducirla. \\
\addlinespace

\textbf{¿Qué pasa si la policía no sigue el procedimiento?} &
La prueba puede quedar excluida o el sospechoso puede ser liberado. Los derechos Miranda deben ser entendidos y comprendidos, no solo enunciados. \\

\midrule
\multicolumn{2}{@{}p{\dimexpr0.94\textwidth+2\tabcolsep\relax}@{}}{%
\textbf{Resumen:} El \textit{U.S. Code} y las \textit{Federal Rules} son federales y orientativos, de modo que cada estado tiene sus propias normas; la Constitución fija un piso mínimo de garantías que los estados pueden ampliar pero no reducir. Las pautas policiales, como los derechos Miranda, exigen que el detenido entienda sus derechos; de lo contrario, la prueba puede excluirse e incluso el sospechoso puede ser liberado.} \\
\end{longtable}

% ============================================================
\chapter[Glosario de términos en inglés]{Glosario de términos en inglés}
% ============================================================

\begin{longtable}{@{}>{\raggedright\arraybackslash}p{0.34\textwidth}>{\raggedright\arraybackslash}p{0.60\textwidth}@{}}
\toprule
\textbf{Término} & \textbf{Significado} \\
\midrule
\endfirsthead
\toprule
\textbf{Término} & \textbf{Significado} \\
\midrule
\endhead
\bottomrule
\endlastfoot

\textit{Grand jury} & Gran jurado: decide si hay indicios para acusar en delitos graves. \\
\textit{Preliminary hearing} & Audiencia preliminar. \\
\textit{Indictment / presentment / information} & Acusación formal (del gran jurado o escrito directo). \\
\textit{Arraignment} & Audiencia en que se informa formalmente la acusación. \\
\textit{Pretrial motion} & Petición previa al juicio (por ejemplo, suprimir prueba). \\
\textit{To suppress evidence} & Excluir pruebas del proceso. \\
\textit{Petit jury} & Jurado ordinario de 12 personas. \\
\textit{Trier of fact} & Quien evalúa los hechos: el jurado. \\
\textit{Opening / closing statement} & Alegato de apertura / de cierre. \\
\textit{Beyond a reasonable doubt} & Más allá de toda duda razonable (estándar penal). \\
\textit{Preponderance of the evidence} & Preponderancia de la prueba (estándar civil). \\
\textit{Hung jury} & Jurado bloqueado, sin veredicto. \\
\textit{Mistrial} & Declaración de nulidad del juicio. \\
\textit{Sentencing} & Determinación de la pena (falso amigo de ``sentencia''). \\
\textit{Aggravating / mitigating factors} & Agravantes / atenuantes. \\
\textit{Suspended sentence} & Pena de ejecución en suspenso. \\
\textit{Restitution} & Reparación económica a la víctima. \\
\textit{Probation} & En Estados Unidos, una pena; en Argentina, suspensión del juicio a prueba. \\
\textit{Allocution} & Última palabra del condenado ante el juez. \\
\textit{Due process of law} & Debido proceso legal. \\
\textit{Double jeopardy} & Prohibición de ser juzgado dos veces por el mismo hecho. \\
\textit{Search / arrest warrant} & Orden judicial de allanamiento / detención. \\
\textit{Shall} & Verbo modal que indica obligación impuesta por una autoridad. \\
\textit{Jury of peers} & Jurado de pares. \\
\end{longtable}

% ============================================================
\chapter[Síntesis general de la clase]{Síntesis general de la clase}
% ============================================================

La clase recorre el proceso penal estadounidense desde las etapas previas hasta la pena y lo contrasta con el proceso civil y con el sistema argentino. Tras el gran jurado y el \textit{indictment}, las \textit{pretrial motions} permiten depurar la prueba, por lo que el respeto de las garantías desde el arresto es decisivo.

El \textit{trial} reproduce la estructura civil (apertura, prueba, cierre), pero con un estándar mucho más exigente (más allá de toda duda razonable) y un jurado de 12 que debe decidir por unanimidad; si no lo logra, el \textit{hung jury} conduce a un \textit{mistrial} y a un nuevo juicio. En delitos gravísimos puede intervenir un tercer jurado para fijar la pena, y el condenado conserva la \textit{allocution} como última palabra.

La segunda parte de la clase vincula estos pasos con el debido proceso, anclado en las Enmiendas Quinta y Sexta, y muestra el carácter federal y orientativo del \textit{U.S. Code} y de las \textit{Federal Rules}, que fijan un piso mínimo de garantías que los estados pueden ampliar pero no reducir, incluidas las pautas policiales como los derechos Miranda. Para quien lee en inglés jurídico, la clave es identificar términos, instituciones y diferencias con el propio sistema.

\end{document}
